% from ega3-1-fr.tex
\oldpage[III]{124}

à un isomorphisme canonique près, puisque
$\mathcal F_k=(i_k)_*(i_k^*(\mathcal F_k))$ (G, II, 4.9.1)~;
l'homomorphisme canonique
$H^n(\widehat X,\widehat{\mathcal F})\to
H^n(X_k,h_k^*(\widehat{\mathcal F}))$ (0, 12.1.3.5) s'écrit donc aussi
\[
H^n(\widehat X,\widehat{\mathcal F})\longrightarrow H^n(X,\mathcal F_k)
\tag{4.1.3.3}\label{III.4.1.3.3-fr}
\]
et ces homomorphismes forment évidemment un système projectif, d'où par
passage à la limite, un homomorphisme canonique
\[
\psi_{n,X}:H^n(\widehat X,\widehat{\mathcal F})
\longrightarrow\varprojlim_kH^n(X,\mathcal F_k).
\tag{4.1.3.4}\label{III.4.1.3.4-fr}
\]
Remplaçant $X$ par un ouvert de la forme $f^{-1}(V)$, où $V$ est un
ouvert affine de $Y$, on a, compte tenu de
(\hyperref[III.1.4.11-fr]{1.4.11}), des homomorphismes
\[
\psi_{n,V}:H^n(\widehat X\cap f^{-1}(V),\widehat{\mathcal F})
\longrightarrow\varprojlim_k\Gamma(V,R^nf_*(\mathcal F_k))~;
\tag{4.1.3.5}\label{III.4.1.3.5-fr}
\]
ces homomorphismes commutent de façon évidente à la restriction de $V$ à
un ouvert affine plus petit, donc définissent finalement un homomorphisme
canonique de faisceaux
\[
\psi_n:R^n\widehat f_*(\widehat{\mathcal F})
\longrightarrow\varprojlim_kR^nf_*(\mathcal F_k).
\tag{4.1.3.6}\label{III.4.1.3.6-fr}
\]
\end{env}

\begin{env}[4.1.4]
\label{III.4.1.4-fr}
Soit enfin $j:\widehat Y\to Y$ le morphisme canonique d'espaces annelés
(I, 10.8.7)~; comme $R^nf_*(\mathcal F)$ est un $\mathcal O_Y$-Module
cohérent (\hyperref[III.3.2.1-fr]{3.2.1}), on a
$j^*(R^nf_*(\mathcal F))=\widehat{R^nf_*(\mathcal F)}$ à un isomorphisme
canonique près (I, 10.8.8) et on a donc un homomorphisme canonique
\[
\rho_n:\widehat{R^nf_*(\mathcal F)}=j^*(R^nf_*(\mathcal F))
\longrightarrow R^n\widehat f_*(i^*(\mathcal F))
=R^n\widehat f_*(\widehat{\mathcal F})
\tag{4.1.4.1}\label{III.4.1.4.1-fr}
\]
défini de façon générale pour les espaces annelés (voir la démonstration
de (\hyperref[III.1.4.15-fr]{1.4.15})). Montrons que le diagramme
\[
\xymatrix{
\widehat{R^nf_*(\mathcal F)}\ar[rr]^{\rho_n}\ar[dr]_{\varphi_n}
&&R^n\widehat f_*(\widehat{\mathcal F})\ar[dl]^{\psi_n}\\
&\displaystyle\varprojlim_kR^nf_*(\mathcal F_k)&
}
\tag{4.1.4.2}\label{III.4.1.4.2-fr}
\]
est commutatif. Il suffit évidemment de prouver la commutativité du
diagramme correspondant d'homomorphismes de préfaisceaux, donc on peut se
borner au cas où $Y$ est affine, et tout revient à prouver que le diagramme
\[
\xymatrix{
\widehat{H^n(X,\mathcal F)}\ar[rr]^{\rho_n}\ar[dr]_{\varphi_n}
&&H^n(\widehat X,\widehat{\mathcal F})\ar[dl]^{\psi_{n,X}}\\
&\displaystyle\varprojlim_kH^n(X,\mathcal F_k)&
}
\tag{4.1.4.3}\label{III.4.1.4.3-fr}
\]
est commutatif. Mais la commutativité de
(\hyperref[III.4.1.3.1-fr]{4.1.3.1}) et les relations vues dans
(\hyperref[III.4.1.3-fr]{4.1.3}) entre les groupes de cohomologie donnent
aussitôt le diagramme commutatif

